\section{Ratios de gestión}
Los ratios de gestión o actividad evalúan la eficiencia del ciclo operativo. Permiten analizar cuánto demora la empresa en cobrar y cuánto tiempo permanece el inventario en almacén.

\subsection{Rotación de cuentas por cobrar y periodo de cobro}
\[
\text{Rotación CxC}=\frac{\text{Ventas netas a crédito}}{\text{Promedio de cuentas por cobrar}}
\]
\[
\text{Periodo promedio de cobro}=\frac{360}{\text{Rotación de cuentas por cobrar}}
\]
\textbf{Caso académico elaborado con fines educativos.} Servicios del Pacífico S.A.C. realizó ventas a crédito por \soles{1,200,000.00}. Sus cuentas por cobrar fueron \soles{180,000.00} al inicio y \soles{220,000.00} al cierre.
\[
\text{Promedio CxC}=\frac{180,000.00+220,000.00}{2}=200,000.00
\]
\[
\text{Rotación}=\frac{1,200,000.00}{200,000.00}=6.00
\quad
\text{Periodo}=\frac{360}{6.00}=60.00\text{ días}
\]
\textbf{Interpretación:} la empresa tarda aproximadamente 60 días en convertir sus ventas a crédito en efectivo.
\textbf{Decisión:} comparar con la política de crédito; si la política es 30 o 45 días, debe reforzarse cobranza.

\newpage
\subsection{Rotación de inventarios}
\[
\text{Rotación de inventarios}=\frac{\text{Costo de ventas}}{\text{Inventario promedio}}
\]
Datos: costo de ventas \soles{900,000.00}, inventario inicial \soles{140,000.00}, inventario final \soles{160,000.00}.
\[
\text{Inventario promedio}=\frac{140,000.00+160,000.00}{2}=150,000.00
\]
\[
\text{Rotación}=\frac{900,000.00}{150,000.00}=6.00
\quad
\text{Permanencia}=\frac{360}{6.00}=60.00\text{ días}
\]
\textbf{Interpretación:} el inventario se renueva seis veces al año y permanece en almacén cerca de 60 días.
